%!TEX root=../../main.tex

\begin{lemma}\label{lemma:lagrange-dual}
    One Lagrange dual of CGL is the following:
    \begin{equation}\label{eq:dual-program}
        \begin{aligned}
            \max_{\eta, \rho} \;
             & - \half (\eta + K \rho - X^\top y) (X^\top X)^+ (\eta + K \rho - X^\top y) + \half \norm{y}_2^2                         \\
             & \quad \quad \text{s.t.} \quad \eta + K \rho \in \Row(X), \; \norm{\eta_\bi}_2 \leq \lambda \; \forall \; \bi \in \calB,
        \end{aligned}
    \end{equation}
    where \( \eta_\bi = \ci - \Ki \ri \) shows that the vectors \( \vi \)
    are, in fact, dual variables.
    Moreover, if \( K = 0 \), then \( \rmin = 0 \) and \( \eta \)
    has the unique solution \( \eta_\bi = \Xbi^\top (y -  Xw) = \ci \).
    That is, the dual parameters are the (unique) block correlation vectors.
\end{lemma}
\begin{proof}
    We re-write the group Lasso problem as follows:
    \begin{align*}
        \min_{w} \half \norm{X z - y}_2^2 + \lambda \sum_{\bi \in \calB} \norm{\wi}_2
        \quad \text{s.t.} \, \, z = w, \, \Ki^\top z_\bi \leq 0.
    \end{align*}
    The Lagrangian for this problem is
    \begin{align*}
        \calL(w, z, \eta, \rho)
         & = \half \norm{X z - y}_2^2 + \abr{\eta, z - w} + \abr{\rho, K^\top z} + \lambda \sum_{\bi \in \calB} \norm{\wi}_2 \\
         & = \half \norm{X z - y}_2^2 + \abr{\eta + K \rho, z} - \abr{\eta, w} + \lambda \sum_{\bi \in \calB} \norm{\wi}_2.
    \end{align*}
    Minimizing over \( z \), we find that stationarity implies
    \[
        X^\top (y - X z) = \eta + K \rho,
    \]
    so that \( \eta + K \rho \in \Row(X) \).
    Solving this system, we find
    \[
        X^\top X z = X^\top y - \eta - K \rho \implies z = (X^\top X)^+ \sbr{X^\top y - \eta - K \rho} + c,
    \]
    where \( c \in \Null(X) \).
    Let us minimize over \( w \) similarly.
    The Lagrangian decouples block-wise in \( w \), so that
    we must solve
    \[
        \min_{\wi} \lambda \norm{\wi}_2 - \abr{\eta_\bi, \wi},
    \]
    for each \( \bi \in \calB \).
    The minimum value is achieved by the (negative) Fenchel conjugate of
    \( \lambda \norm{\wi}_2 \)
    evaluated at \( \eta_\bi \);
    that is,
    \[
        \min_{\wi} \lambda \norm{\wi}_2 - \abr{\eta_\bi, \wi} = - \mathds{1}(\norm{\eta_\bi}_2 \leq \lambda).
    \]
    Combining this with the expression for \( z \), we obtain
    \begin{align*}
        \calL(c, \eta, \rho)
         & = - \half (\eta + K \rho - X^\top y) (X^\top X)^+ (\eta + K \rho - X^\top y)
        + \abr{\eta + K \rho, c}
        \\
         & \hspace{4em} - \sum_{\bi \in \calB} \mathds{1}(\norm{\eta_\bi}_2 \leq \lambda) \\
         & = - \half (\eta + K \rho - X^\top y) (X^\top X)^+ (\eta + K \rho - X^\top y)
        - \sum_{\bi \in \calB} \mathds{1}(\norm{\eta_\bi}_2 \leq \lambda),
    \end{align*}
    where the second equality follows since \( \eta + K \rho \in \Row(X) \) and \( c \in \Null(X) \)
    are orthogonal.
    (Alternatively, one can observe that the dual problem is unbounded below
    whenever \( \abr{c, \eta + K \rho} \neq 0 \).)
    Thus, the dual problem is equal to
    \begin{align*}
         & \max_{\eta, \rho} - \half (\eta + K \rho - X^\top y) (X^\top X)^+ (\eta _ K \rho - X^\top y)
        - \sum_{\bi \in \calB} \mathds{1}(\norm{\eta_\bi}_2 \leq \lambda)                               \\
         & \hspace{4em} - \mathds{1}(\eta + K \rho \in \Row(X)),
    \end{align*}
    which completes the derivation.

    Recalling \( z = w \) for any primal-dual optimal pair and
    \[
        X^\top (y - X z) = \eta + K \rho,
    \]
    shows that \( \eta_\bi = \ci - K \rho \) as claimed.
    Moreover, if \( K = 0 \), then we may assume without loss of generality
    that he corresponding dual vectors \( \ri \) are zero.
    In this case, \( \eta_\bi = \ci \) and,
    since \( \ci \) is unique, the dual solution must also be unique.
\end{proof}

\begin{proposition}\label{prop:non-negative-regression}
    Let \( \lambda > 0 \) and \( w \in \solfn(\lambda) \).
    If \( \wi = 0 \), then any solution to \cref{eq:non-negative-regression}
    is dual optimal for block \( \bi \).
\end{proposition}
\begin{proof}
    Let \( \ri \) be a solution to \cref{eq:non-negative-regression}.
    Since \( w \) is optimal and strong duality holds, there exists
    some min-norm dual optimal vector \( \rmin \).
    Moreover \( \rmin \) satisfies \( \rmin_\bi \geq 0 \) and
    \[
        \norm{\Ki \ri - \ci}^2_2 \leq \norm{\Ki \rmin_\bi - \ci}^2_2 \leq \lambda^2,
    \]
    so \( w \) is both feasible satisfies stationarity of the Lagrangian,
    Finally, because \( \wi = 0 \), complementary slackness,
    \[
        [\ri]_j \cdot [\Ki]_j^\top \wi = 0,
    \]
    is verified for every \( j \in [a_\bi] \).
    Since the KKT conditions are sufficient for primal-dual optimality, we
    conclude that \( \ri \) is dual optimal.
    This completes the proof.
\end{proof}

